\section{智能眼镜为什么可能成为入口}

上一章说明 Rokid 为什么离开音箱，本节解释它为什么押注眼镜。祝铭明认为，AI 真正有价值的使用方式是随时随地、碎片化、always on。手机虽然强，但仍有一套解锁、打开 app、进入功能的 GUI 路径；智能眼镜如果一直戴着，就可以把许多短任务变成“所要即所得”：看时间、收信息、扫码、翻译、识别对象、查询背景，都不需要掏手机。

这不是说手机会马上消失，而是交互频率会重分配。手机可能保留通信、重任务、复杂输入和大屏操作；眼镜会先接管大量五分钟以内的碎片交互。祝铭明给出的判断是：对已经戴眼镜的人群而言，眼镜加 AI 可能在三到五年内成为普及性东西；对不愿意戴眼镜的人，不要强行说服，而是等能力差距足够明显后再由用户选择。

\begin{figure}[H]
\centering
\includegraphics[width=0.96\textwidth]{smart-glasses-stage.png}
\caption{智能眼镜阶段判断：从黑莓到 iPhone 1 之间，入口正在形成。自制概念图，依据 01:05:45--01:13:05 对谈内容整理。}
\end{figure}

\begin{knowledgebox}{读图：阶段不成熟，不等于方向不成立}
图中从 Feature 到 BlackBerry、iPhone 1、Platform、Mass Market，表达的是智能眼镜还在早期入口形成阶段。今天的产品未必成熟，但如果使用时长、开发者和碎片任务迁移成立，它就有机会从功能硬件走向平台。
\end{knowledgebox}

\subsection{信息直达：把系统成本压到最低}

本节把“方便”拆成可分析的交互成本。祝铭明举了扫码、查天气、看时间和看消息的例子：手机路径通常包括拿出设备、解锁、找到 app、进入功能、确认操作；眼镜如果一直佩戴，就可以直接语音唤起、视觉确认或轻触查看。几秒差异在单次任务中不大，但当任务都是五分钟以内碎片动作时，几秒会累积成入口迁移。

\[
T_{\text{task}} = T_{\text{wake}} + T_{\text{path}} + T_{\text{input}} + T_{\text{output}} + T_{\text{confirm}}
\]

其中，$T_{\text{task}}$ 表示完成一个短任务的总时间；$T_{\text{wake}}$ 是拿出或唤醒设备的时间；$T_{\text{path}}$ 是从系统入口走到具体功能的路径成本；$T_{\text{input}}$ 是表达意图的时间；$T_{\text{output}}$ 是接收结果的时间；$T_{\text{confirm}}$ 是确认和纠错的时间。眼镜路线的核心，不是让每一项都归零，而是把 $T_{\text{wake}}$ 和 $T_{\text{path}}$ 压到很低，并用显示降低 $T_{\text{output}}$ 和 $T_{\text{confirm}}$ 的不确定性。

\begin{importantbox}{入口不是单次效率，而是频率乘以效率}
如果一个动作每天只发生一次，节省五秒没有意义；如果通知、时间、扫码、翻译、识别、搜索每天发生几十次，路径成本就会变成入口选择。Rokid 的判断是：智能眼镜先替代大量碎片化轻任务，而不是立刻替代手机的重任务。
\end{importantbox}

\subsection{显示为什么重要}

中美智能眼镜产品定义的核心差异之一，是是否需要显示。Rokid 认为，如果眼镜是 AI 入口，显示是必选项。以翻译为例，没有显示时，系统必须等到有足够信心再把翻译读出来，交流中会出现几秒到十秒空档；有显示时，可以边听边显示、边修正，用户能容忍早期不完全准确的中间结果。显示不是为了炫技，而是为了让 AI 输出从纯语音变成可扫描、可纠错、可并行的界面。

\begin{knowledgebox}{术语消化：AUI 与 GUI}
\noindent\begin{tabularx}{\textwidth}{p{0.18\textwidth}p{0.34\textwidth}X}
\toprule
交互方式 & 优势 & 限制 \\
\midrule
GUI & 可见、可编辑、可扫描，适合复杂任务 & 需要路径和界面操作，系统成本高。 \\
AUI & 语音直达，适合短任务和免手操作 & 输出串行，不适合长内容和低置信度中间结果。 \\
眼镜显示 + AI & 语音输入与视觉输出结合 & 需要佩戴舒适、续航、显示和隐私共同过线。 \\
\bottomrule
\end{tabularx}
\end{knowledgebox}

\begin{warningbox}{不要把显示理解成“多一个屏幕”}
显示的真正价值不是把手机屏幕搬到眼前，而是改变 AI 输出的时间结构。语音输出必须串行、完整、较高置信度；视觉输出可以逐步显示、允许中间修正、让用户扫读。翻译场景之所以典型，是因为它把延迟、置信度和可纠错性同时暴露出来。
\end{warningbox}

\noindent\begin{tabularx}{\textwidth}{p{0.18\textwidth}p{0.34\textwidth}X}
\toprule
场景 & 无显示路径 & 有显示路径 \\
\midrule
翻译 & 等模型有足够信心后再读出，双方对话中断 & 可先显示粗译，再随上下文修正，用户可扫读。 \\
通知 & 语音播报会打扰场景，手机查看又增加路径 & 眼前短暂提示，用户可忽略或进一步处理。 \\
扫码/支付 & 需要打开手机、找入口、确认 & 视线对准二维码，语音或轻触确认。 \\
识别对象 & 手机拍照或打开相机入口 & 眼镜直接利用视线方向形成查询。 \\
\bottomrule
\end{tabularx}

\subsection{佩戴边界：不要说服不想戴眼镜的人}

本节处理一个现实反对意见：很多人不愿意戴眼镜，甚至近视也选择隐形或手术。祝铭明的答案很克制：当前阶段不要强行说服这批人。先服务已经接受有框眼镜的人群，让他们在原有佩戴行为上叠加 AI 能力；当身边戴眼镜的人获得明显信息优势时，再让不愿戴的人重新选择。这里的产品策略不是“全民立刻佩戴”，而是从更自然的原住民人群开始。

\begin{knowledgebox}{实践经验：采用曲线从最少阻力人群开始}
硬件产品的早期用户不一定是最大人群，而是阻力最小、反馈最强、愿意忍受早期不完美的人群。对智能眼镜来说，愿意戴有框眼镜的人、科技极客、出差/翻译/会议高频用户，比完全抗拒眼镜的人更适合作为早期主用户。
\end{knowledgebox}

\noindent\begin{tabularx}{\textwidth}{p{0.20\textwidth}p{0.34\textwidth}X}
\toprule
用户类型 & 初期阻力 & 产品策略 \\
\midrule
本来戴眼镜的人 & 最低，只需把原有佩戴升级为 AI 能力 & 先做舒适、度数、重量和日常使用。 \\
愿意部分场景佩戴的人 & 中等，需要关键场景足够强 & 旅行、翻译、会议、导航、扫码等高价值场景。 \\
不愿戴眼镜的人 & 最高，生理和心理阻力都存在 & 不强行教育，等能力差距和形态成熟后再扩展。 \\
\bottomrule
\end{tabularx}

\subsection{中美产品定义差异}

美国很多智能眼镜更像 sunglasses 逻辑：户外、社交、拍照、视频分享，不一定 always on，也不一定把显示作为必选项。Rokid 更偏中国眼镜逻辑：大量用户本来就戴眼镜，室内外长时间佩戴更自然；因此产品定义从第一天就偏 AI 加 AR，强调显示、助手和随身信息入口。Meta/Ray-Ban 的路线是社交内容生产工具，Rokid 的路线是个人信息消费和 AI 助手。

\begin{figure}[H]
\centering
\includegraphics[width=0.94\textwidth]{china-us-glasses-definition.png}
\caption{中美定义智能眼镜的差异：产品定义、渠道和用户场景不同。自制概念图，依据 01:23:38--01:41:35 对谈内容整理。}
\end{figure}

\begin{knowledgebox}{读图：同样长得像眼镜，不代表是同一物种}
左侧中国路线强调近视人群、全天佩戴、显示和 AI 助手；右侧美国路线更容易从 sunglasses、社交拍摄和品牌时尚进入。两者可能最终交汇，但早期产品定义、用户场景和功能优先级不同。
\end{knowledgebox}

\begin{importantbox}{产品定义决定技术栈}
如果产品是社交拍摄工具，摄像头、品牌、分享链路和隐私提示是核心；如果产品是 AI 入口，显示、OS、低延迟反馈、功耗、佩戴舒适和模型接入是核心。同样叫智能眼镜，技术优先级可能完全不同。
\end{importantbox}

\subsection{本章小结}

智能眼镜的入口逻辑是 always-on、信息直达和碎片任务迁移。显示能力决定 AI 输出是否能从串行语音变成可视、可纠错、可平台化的交互；佩戴边界则决定早期人群不能被无限扩大，必须先从阻力最小的用户开始。

